\chapter{Ingesta y fusión de fuentes heterogéneas}
\label{cap:ingesta}

El presente capítulo describe cómo se adquieren, generan y aterrizan las
dos fuentes de datos del sistema. Se recorren la doble adquisición de
datos de ORCID, la generación de currículos CVN sintéticos sembrados con
esos mismos datos, el aterrizaje en bronze con procedencia y la
verificación de extremo a extremo del flujo de Airflow que encadena las
cuatro tareas. El enlace deliberado entre el CVN sintético y su semilla
ORCID, introducido en este capítulo, es la clave de unión que la
resolución de identidad del Capítulo~\ref{cap:procesamiento} explota.

\section{Arquitectura de la ingesta}
\label{sec:ingesta_arquitectura}

La ingesta se implementa como un único flujo de Airflow,
\texttt{ingest\_validate}, con cuatro tareas encadenadas: obtención del
subconjunto ORCID masivo, generación de CVN sintético, enriquecimiento vía
la API de ORCID y validación con aterrizaje en bronze. Este orden no es
arbitrario. La tarea de enriquecimiento consulta la API por los
identificadores ORCID que los CVN sintéticos ya declaran, por lo que no
puede ejecutarse antes de que existan. Ese acoplamiento deliberado entre tareas
es la propia demostración de fusión dirigida entre fuentes: el
identificador que une un CVN sintético con su semilla real es el mismo que
la tarea de enriquecimiento usa para pedir datos frescos a la API. La
Figura~\ref{fig:tfm_ingest_dag} resume el orden de las cuatro tareas.

\begin{figure}[H]
  \centering
  \begin{tikzpicture}[
    font=\small,
    node distance=7mm,
    task/.style={rectangle, draw=azul, thick, rounded corners, fill=gris95,
      text width=7.4cm, align=center, inner sep=5pt},
    >=Stealth
  ]
    \node[task] (t1) {\texttt{fetch\_orcid\_bulk\_subset}\\
      obtención del subconjunto ORCID masivo};
    \node[task, below=of t1] (t2) {\texttt{generate\_synthetic\_cvn}\\
      generación de CVN sintético};
    \node[task, below=18mm of t2] (t3) {\texttt{fetch\_orcid\_api\_enrichment}\\
      enriquecimiento vía la API de ORCID};
    \node[task, below=of t3] (t4) {\texttt{validate\_and\_land\_bronze}\\
      validación y aterrizaje en bronze};

    \draw[->] (t1) -- (t2);
    \draw[->] (t2) -- node[right, text width=4cm, align=left, font=\scriptsize]
      {usa los identificadores ORCID que los CVN sintéticos ya declaran}
      (t3);
    \draw[->] (t3) -- (t4);
  \end{tikzpicture}
  \caption{Orden de tareas del flujo \texttt{ingest\_validate}.}
  \label{fig:tfm_ingest_dag}
\end{figure}

Cada tarea corre en un pod independiente lanzado con
\texttt{KubernetesPodOperator}, sobre una imagen propia con Python 3.14 y
las dependencias del proyecto. El propio Airflow corre sobre Python 3.13,
así que la lógica de negocio del proyecto no puede ejecutarse en sus
\textit{workers} y necesita, igual que la ejecución de Spark del
Capítulo~\ref{cap:infraestructura}, un pod dedicado por tarea. El código
del proyecto, los esquemas y los datos llegan a cada pod montados
directamente desde el disco del nodo, sin necesidad de reconstruir la
imagen en cada cambio de código. Las tareas se comunican entre sí a través
de ese mismo directorio de datos compartido, con un subdirectorio por
ejecución. El mecanismo nativo de paso de mensajes de Airflow se reserva
para resúmenes pequeños, como recuentos y rutas, y nunca para los propios
datos, que alcanzan varios gigabytes por ejecución.

\section{Adquisición de datos de ORCID}
\label{sec:ingesta_orcid}

ORCID~\cite{orcid_about} ofrece dos mecanismos públicos y oficiales para
obtener datos, y el sistema usa ambos porque cubren necesidades
distintas y complementarias. Ninguno de los dos sustituye al otro.

\begin{itemize}
  \item \textbf{API pública de ORCID~\cite{orcid_public_api_docs}.}
  Resuelve consultas puntuales por identificador y da acceso al registro
  más reciente de una persona concreta. Su límite de tasa, en el nivel
  anónimo sin autenticación, es de 25\,000 lecturas al día y 12 peticiones
  por segundo por dirección IP. Durante la implementación se encontró una
  discrepancia entre la documentación pública sobre el registro de
  aplicaciones cliente, que describe el nivel autenticado como accesible a
  cualquier investigador individual, y el formulario de registro real, que
  exige describir la aplicación como herramienta de una organización
  miembro. Ese requisito no encaja con un proyecto académico personal, de
  modo que el sistema usa el nivel anónimo, sin que ello afecte a la
  demostración de fusión dirigida entre fuentes.
  \item \textbf{Fichero público de datos de ORCID.} ORCID publica
  anualmente un volcado completo de sus perfiles públicos en la
  plataforma de datos abiertos Figshare~\cite{figshare_docs}. Este
  fichero da volumen real en lugar de tener que sintetizar también los
  datos de ORCID. El fichero completo se distribuye en doce archivos, de
  los cuales solo se usa el de resúmenes de perfil, de unos 46 GB, y no
  los once de actividades detalladas, de unos 191 GB en conjunto, que
  exceden lo que este trabajo necesita. El sistema recorre el fichero de
  resúmenes completo, sin detenerse al alcanzar un volumen objetivo, para
  no introducir un sesgo de muestreo hacia los primeros registros del
  archivo, y filtra los perfiles con afiliación declarada en
  España~\cite{orcid_public_data_file_2025}. La integridad del fichero
  descargado se verifica contra la suma de comprobación que la propia
  Figshare publica junto a él.
\end{itemize}

La Tabla~\ref{tab:tfm_mecanismos_orcid} resume los dos mecanismos.

\begin{table}[H]
  \centering
  \small
  \renewcommand{\arraystretch}{1.2}
  \begin{tabular}{
    p{0.18\textwidth}
    >{\raggedright\arraybackslash}p{0.28\textwidth}
    >{\raggedright\arraybackslash}p{0.22\textwidth}
    >{\raggedright\arraybackslash}p{0.18\textwidth}}
    \hline
    \textbf{Mecanismo} & \textbf{Uso} & \textbf{Volumen} & \textbf{Límite} \\
    \hline
    API pública & Enriquecimiento dirigido por identificador &
      Un registro por consulta &
      25\,000 lecturas/día, 12 peticiones/s \\
    Fichero de datos & Volumen real de siembra para la generación sintética &
      301\,763 perfiles con afiliación española filtrados &
      Instantánea anual, sin límite de tasa \\
    \hline
  \end{tabular}
  \caption{Los dos mecanismos públicos de adquisición de datos de ORCID.}
  \label{tab:tfm_mecanismos_orcid}
\end{table}

El filtrado del fichero completo, sobre 26\,078\,951 registros escaneados,
dejó 301\,763 perfiles con afiliación española, un 1,157\% del total. La
primera implementación, que procesaba el archivo en flujo directamente
desde la red, proyectaba unas siete horas de ejecución. El cuello de
botella no estaba en la red sino en el análisis de XML de cada registro,
así que la solución fue descargar antes el archivo al disco local y
aplicar un filtro previo a nivel de bytes que descarta un registro sin
analizarlo completamente cuando resulta evidente que no puede coincidir.
Con ese cambio, la extracción completa tardó 81,6 minutos. El resultado se
guarda como un fichero XML por perfil, organizado en subdirectorios por
prefijo del identificador, que sirve de semillero a la generación de CVN
sintéticos descrita a continuación.

\section{Generación de currículos CVN sintéticos}
\label{sec:ingesta_sintetico}

El sistema necesita un volumen realista de currículos CVN para demostrar
la plataforma a escala, y ese volumen no puede provenir de currículos
reales por las razones de privacidad expuestas en el
Capítulo~\ref{cap:introduccion}. La solución adoptada genera documentos
Open CVN sintéticos, válidos contra el mismo esquema JSON
Schema~\cite{json_schema_org} que define el formato oficial del proyecto,
y los siembra con campos públicos reales tomados del subconjunto ORCID
descrito en la Sección~\ref{sec:ingesta_orcid}: nombre, afiliaciones,
títulos de obras y palabras clave. Nunca se copian de ORCID la biografía,
el correo electrónico ni las URL de investigador de un perfil, porque
ORCID pide explícitamente no presentar datos modificados como genuinos,
aunque su fichero de datos se publique bajo licencia CC0. El resto de la
identidad de cada documento, como el teléfono o la fecha del documento,
se inventa por completo, con valores que no pueden confundirse con datos
reales, y cada documento generado queda marcado como sintético en sus
propios metadatos.

Una parte configurable de los documentos generados, un 70\% en la
configuración de referencia, declara el identificador ORCID de su
semilla en el campo del esquema reservado a identificadores digitales de
autor, con el código de fuente que las tablas de referencia del CVN
asignan a ORCID. El resto conserva el nombre y las afiliaciones de su
semilla, con variaciones realistas, pero sin ese identificador. Un fichero
de control aparte, no aterrizado en el almacenamiento de objetos, registra
para cada documento generado el identificador de su semilla y el tipo de
enlace aplicado. Esa combinación de documentos con enlace declarado,
documentos con coincidencia solo por nombre y afiliación, y documentos sin
relación entre sí, es precisamente la variedad de casos que la resolución
de identidad del Capítulo~\ref{cap:procesamiento} necesita para
verificarse, y el fichero de control es su verdad de referencia. Cuando el
volumen solicitado supera el número de semillas disponibles, una semilla
se reutiliza para varios documentos, con una selección distinta de sus
obras y una identidad inventada distinta en cada uno, pero sin alterar
fechas, organizaciones ni otros hechos tomados de ORCID.

La validación de cada documento generado ocurre en dos capas. La primera
reutiliza el validador Open CVN del TFG~\cite{tfg_open_cvn} sobre el
documento completo. La segunda, propia de este generador, valida además
cada entrada contra la definición estricta de su propio tipo de entidad
dentro del esquema. Esta segunda capa resultó necesaria en la práctica y
no solo en el diseño. El esquema del documento completo deja el contenido
de cada entrada sin restricciones de campos, por lo que un documento con
campos inventados pasaría la primera capa sin problema. La primera
ejecución de prueba generó 200 documentos con solo dos errores de
construcción, y aun así 179 de ellos resultaron inválidos frente a la
segunda capa, por dos campos con nombre incorrecto y una fecha vacía en un
campo obligatorio. Corregidos ambos errores, las ejecuciones posteriores
no descartaron ningún documento. Una ejecución de 10\,000 documentos, con
semilla fija y el mismo reparto de enlace del 70\%, generó 7\,000
documentos enlazados y 3\,000 sin enlazar, sin ningún documento inválido,
en 169 segundos.

\section{Aterrizaje en bronze con procedencia}
\label{sec:ingesta_bronze}

Las tres fuentes de datos, el subconjunto ORCID masivo, el enriquecimiento
por API y los CVN sintéticos, aterrizan como objetos de texto plano en
formato JSON Lines bajo el prefijo \texttt{bronze} del cubo de MinIO
configurado en el Capítulo~\ref{cap:infraestructura}, organizados por
fuente y por fecha de ingesta. No se escriben como tablas Iceberg en esta
fase: mantener el aterrizaje crudo separado del catálogo de tablas evita
que la lógica interna de Iceberg tropiece con objetos ajenos a sus propias
tablas, y deja la creación de las tablas bronze de Iceberg como parte del
trabajo de transformación del Capítulo~\ref{cap:procesamiento}.

Cada registro aterrizado se envuelve en un sobre de procedencia común a
las tres fuentes, resumido en la Tabla~\ref{tab:tfm_sobre_procedencia}. El
propio contenido original, ya sea el XML de ORCID sin modificar, el JSON
que devuelve la API o el documento CVN sintético, viaja como carga útil
dentro de ese sobre. Este vocabulario de procedencia extiende el mismo
utilizado por el bloque de trazabilidad del formato Open CVN del TFG, en
lugar de inventar una convención nueva.

\begin{table}[H]
  \centering
  \small
  \renewcommand{\arraystretch}{1.2}
  \begin{tabular}{
    p{0.28\textwidth}
    >{\raggedright\arraybackslash}p{0.58\textwidth}}
    \hline
    \textbf{Campo} & \textbf{Contenido} \\
    \hline
    Identificador de registro & Identifica el registro de forma única dentro de su fuente \\
    Fuente de origen & Cuál de las tres fuentes produjo el registro \\
    Referencia al origen & Referencia concreta dentro de esa fuente, por ejemplo un identificador ORCID \\
    Momento de obtención & Cuándo se obtuvo el dato de su fuente original \\
    Momento de aterrizaje & Cuándo se aterrizó el registro en bronze \\
    Identificador de ejecución & Qué ejecución del flujo aterrizó el registro \\
    Resultado de la comprobación & Estado de la validación descrita a continuación \\
    \hline
  \end{tabular}
  \caption{Campos del sobre de procedencia, comunes a las tres fuentes.}
  \label{tab:tfm_sobre_procedencia}
\end{table}

Antes de aterrizar, cada registro pasa una comprobación de calidad ligera,
deliberadamente más simple que la resolución de identidad del
Capítulo~\ref{cap:procesamiento}. Para ORCID exige un XML bien formado, un
identificador con dígito de control válido y al menos un nombre. Para el
CVN sintético exige que el documento pase el validador Open CVN del TFG.
Los registros que no superan esta comprobación se separan a una zona de
rechazados, con sus errores adjuntos, en lugar de descartarse sin dejar
rastro, y la propia tarea de aterrizaje falla si la proporción de
rechazados de una ejecución supera un umbral del 5\%, para convertir un
fallo sistemático de origen en un error visible en vez de en un aterrizaje
silenciosamente incompleto. Sobre los primeros 20\,000 registros ORCID
reales, la tasa de rechazo medida fue del 2,66\%, y su causa casi
exclusiva fue la ausencia de un apellido público en el perfil.

Aterrizar el fichero masivo completo cada día no tiene sentido: su
contenido no cambia entre ejecuciones y ocupa varias decenas de gigabytes.
Por ello, el subconjunto masivo solo se aterriza una vez por instantánea
del fichero de origen, con un límite de 20\,000 registros por defecto que
mantiene cada ejecución diaria en un tamaño manejable. Ese límite es un
parámetro de la propia ejecución, así que puede elevarse para obtener un
volumen mayor cuando haga falta, como en el trabajo de rendimiento del
Capítulo~\ref{cap:visualizacion}.

\section{Verificación de extremo a extremo}
\label{sec:ingesta_verificacion}

Igual que en el Capítulo~\ref{cap:infraestructura}, la verificación del
flujo completo se apoya en dos comprobaciones independientes. La primera
es el propio resultado de cada tarea del flujo de Airflow, con sus
recuentos de registros generados, aceptados y rechazados. La segunda es
una lectura de los mismos datos desde un proceso Spark ajeno a la
ejecución que los produjo, que debe coincidir exactamente con esos
recuentos.

Una ejecución con los parámetros por defecto aterrizó 19\,469 registros
del subconjunto ORCID masivo en cuarenta fragmentos, junto con 531
registros rechazados, 1\,000 documentos CVN sintéticos y 200 respuestas de
la API de enriquecimiento, con todos los registros aterrizados portando
los doce campos de su sobre de procedencia y sin ningún identificador de
registro duplicado. Una lectura independiente con Spark de esos mismos
objetos, ejecutada desde un pod distinto, contó exactamente esas mismas
cifras. Esa misma lectura reveló un problema real de diseño: Spark lee
todos los objetos bajo un prefijo, tenga o no fichero de manifiesto, así
que una ejecución que había fallado el umbral de rechazo dejaba sus
registros ya aterrizados visibles para cualquier lector que ignorase el
manifiesto. La regla de que un fichero de manifiesto ausente señala una
partición incompleta era, hasta entonces, solo una convención. La corrección
aplicada borra los objetos válidos ya aterrizados de una ejecución que
falla. La zona de rechazados de esa misma ejecución no se toca, y sus
registros descartados quedan disponibles para diagnóstico igual que en
cualquier ejecución que no falla. Una repetición de la misma prueba, tras
el cambio, no dejó ningún objeto huérfano en el almacenamiento.

Esta forma de trabajar deja una limitación conocida y aceptada para este
alcance. La semilla de generación sintética se mantiene fija entre
ejecuciones para dar una base de comparación reproducible, así que cada
ejecución diaria del flujo aterriza los mismos documentos sintéticos en
una nueva partición por fecha, y corresponde a la transformación del
Capítulo~\ref{cap:procesamiento} deduplicar por identificador de registro
o leer solo la partición más reciente.
